% -*- coding: utf-8 -*-
\chapter{面向 6G 的 RIS 主被动波束赋形与信道估计机制}
\label{chap:methodology}

针对第二章所揭示的在 6G 高频多用户场景下无源可重构智能表面（RIS）面临的“主被动波束非凸高维耦合”与“无射频链级联信道获取开销爆炸”两大核心理论与工程挑战，本章系统性构建了一套协同优化的物理层传输理论体系。本章的总体思路可以概括为「先建模、再解耦、后估计、末论证」四步：建模给出可分析的数学对象，解耦把不可解的双变量非凸问题拆成可解的凸子问题，估计解决「这些信道参数从哪来」的前置问题，论证则回答「迭代是否会停、估计精度还有多少余量、算力是否够用」三个可实施性问题。

具体而言，本章首先建立严谨的 BS-RIS-User 多用户 MISO 级联信道与信号传输数学模型，明确直连链路与级联反射链路的叠加结构，并指出级联矩阵的秩一哈达玛调制特性——这一特性既是后续算法可以利用的结构，也是信道估计稀疏先验失效的根源；接着详细阐述基于交替优化（AO）与二次变换分式规划（Fractional Programming, FP）的主被动波束协同解耦优化算法，将基站预编码与 RIS 相移分离为两个可交替求解的子问题，其中基站侧得到闭式最优解、RIS 侧在复圆流形上做共轭梯度下降；随后提出基于离散傅里叶变换（DFT）正交反射码本的低开销级联信道估计协议，把导频开销压到 $P = N+1$ 的理论下界，并说明在部分开销下压缩感知先验为何不天然成立；最后推导信道估计克拉美罗下界（CRLB），并在非凸优化理论框架下证明交替优化算法的单调收敛性与计算复杂度界限。

为保证推导的可复核性，本章的符号约定如下：$\bm{G}\in\mathbb{C}^{N\times M}$ 表示基站至 RIS 的链路矩阵，$\bm{h}_{r,k}\in\mathbb{C}^{N}$ 表示 RIS 至用户 $k$ 的链路向量，$\bm{h}_{d,k}\in\mathbb{C}^{M}$ 为基站至用户 $k$ 的直连向量，$\bm{v}\in\mathbb{C}^{N}$ 为 RIS 反射相移向量且满足恒模约束 $|v_n| = 1$，$\bm{w}_k\in\mathbb{C}^{M}$ 为基站对用户 $k$ 的预编码向量，$P_{\max}$ 为总功率上限，$\sigma^2$ 为噪声方差。全文的「和速率」均指以 bps/Hz 为单位的可达速率之和 $R_{\mathrm{sum}} = \sum_k \log_2(1+\gamma_k)$，与第四章实验中的单位保持一致。

% =============================================================================
% 模块化分块深写小节引入 (Modular Deep Subsections)
% =============================================================================
% -*- coding: utf-8 -*-
\section{系统模型与问题形式化表征}
\label{sec:system_model}

考虑如图所示的典型下行多用户单小区多天线无线通信系统。基站（BS）部署有 $M$ 根均匀线阵（ULA）全向发射天线，服务区域内随机分布着 $K$ 个单天线移动通信用户，记为集合 $\mathcal{K} = \{1, 2, \dots, K\}$。由于建筑物密集体、植被遮挡或人体阻隔，基站至部分移动用户之间的视距（LoS）直接链路可能发生严重衰落或完全阻断。为此，在基站视线良好的建筑物外立面部署一块配备 $N$ 个亚波长无源反射单元的可重构智能表面（RIS），通过自适应协同重构空间电磁波反射路径，为主覆盖盲区用户提供高品质虚拟直射通道。

\subsection{传输信号与信道传播模型}
基站向 $K$ 个用户并发传输下行独立数据符号。令 $s_k \in \mathbb{C}$ 表示基站发送给第 $k$ 个用户的归一化信息数据符号，满足统计独立假设：
\begin{equation}
  \mathbb{E}[|s_k|^2] = 1, \quad \mathbb{E}[s_k s_j^*] = 0, \quad \forall k \neq j
  \label{eq:data_symbol_property}
\end{equation}
基站端采用线性数字预编码技术，令 $\bm{w}_k \in \mathbb{C}^{M \times 1}$ 表示面向用户 $k$ 的有源发射波束赋形向量，则基站发射的时域复基带信号向量 $\bm{x} \in \mathbb{C}^{M \times 1}$ 可表征为：
\begin{equation}
  \bm{x} = \sum_{k=1}^K \bm{w}_k s_k = \bm{W} \bm{s}
  \label{eq:bs_transmit_signal}
\end{equation}
其中 $\bm{W} = [\bm{w}_1, \bm{w}_2, \dots, \bm{w}_K] \in \mathbb{C}^{M \times K}$ 为基站联合发射预编码矩阵，$\bm{s} = [s_1, s_2, \dots, s_K]^T \in \mathbb{C}^{K \times 1}$ 为多用户数据向量。受限于基站射频功率放大器（HPA）的线性动态范围与能耗预算，发射信号必须满足最大总发射功率约束：
\begin{equation}
  \mathbb{E}[\|\bm{x}\|^2] = \mathrm{Tr}(\bm{W}\bm{W}^H) = \sum_{k=1}^K \|\bm{w}_k\|^2 \le P_{\max}
  \label{eq:power_constraint}
\end{equation}
其中 $P_{\max} > 0$ 为基站额定发射功率上界。

令 $\bm{h}_{d,k}^H \in \mathbb{C}^{1 \times M}$ 表示基站至用户 $k$ 的直接下行信道向量；$\bm{G} \in \mathbb{C}^{N \times M}$ 表示基站至 RIS 的多输入多输出无线信道矩阵；$\bm{h}_{r,k}^H \in \mathbb{C}^{1 \times N}$ 表示 RIS 至用户 $k$ 的反射无线信道响应向量。各反射单元对入射信号施加独立的相移与幅度衰减，RIS 整体电磁反射响应由对角矩阵 $\bm{\Theta} \in \mathbb{C}^{N \times N}$ 描述：
\begin{equation}
  \bm{\Theta} = \mathrm{diag}\left( \beta_1 e^{j\theta_1}, \beta_2 e^{j\theta_2}, \dots, \beta_N e^{j\theta_N} \right)
  \label{eq:ris_theta_matrix}
\end{equation}
其中 $\theta_n \in [0, 2\pi)$ 表示第 $n$ 个反射单元的相移控制角，$\beta_n \in [0, 1]$ 为其幅度反射系数。在无源理想超表面设计中，为实现入射电磁能量的最大化无损反射，通常设置理想全反射状态，即 $\beta_n = 1, \forall n=1,\dots,N$。

\subsection{多用户接收信号与复合级联信道表示}
经由直连路径与 RIS 反射路径叠加后，第 $k$ 个移动用户处的接收信号 $y_k \in \mathbb{C}$ 可建模为：
\begin{equation}
  y_k = \left( \bm{h}_{d,k}^H + \bm{h}_{r,k}^H \bm{\Theta} \bm{G} \right) \bm{x} + n_k
  \label{eq:user_received_signal}
\end{equation}
其中 $n_k \sim \mathcal{CN}(0, \sigma_k^2)$ 为用户 $k$ 接收机处服从均值为零、方差为 $\sigma_k^2$ 的加性独立复高斯白噪声。

定义 RIS 的无源等效反射波束向量为 $\bm{v} = [v_1, v_2, \dots, v_N]^H \in \mathbb{C}^{1 \times N}$，其中 $v_n = \beta_n e^{-j\theta_n}$，则有：
\begin{equation}
  \bm{h}_{r,k}^H \bm{\Theta} \bm{G} = \bm{v}^H \mathrm{diag}(\bm{h}_{r,k}^H) \bm{G}
  \label{eq:cascaded_algebraic_identity}
\end{equation}
定义基站经 RIS 反射至用户 $k$ 的复合级联信道矩阵为：
\begin{equation}
  \bm{H}_{\mathrm{casc},k} \triangleq \mathrm{diag}(\bm{h}_{r,k}^H) \bm{G} \in \mathbb{C}^{N \times M}
  \label{eq:cascaded_channel_def}
\end{equation}
则基站到用户 $k$ 的全等效复合信道向量 $\bm{h}_k^H(\bm{v}) \in \mathbb{C}^{1 \times M}$ 可紧凑表示为关于反射向量 $\bm{v}$ 的线性函数：
\begin{equation}
  \bm{h}_k^H(\bm{v}) \triangleq \bm{h}_{d,k}^H + \bm{v}^H \bm{H}_{\mathrm{casc},k}
  \label{eq:equivalent_composite_channel}
\end{equation}

将式~\eqref{eq:bs_transmit_signal}与式~\eqref{eq:equivalent_composite_channel}代入式~\eqref{eq:user_received_signal}，用户 $k$ 的接收信号进一步展开为：
\begin{equation}
  y_k = \bm{h}_k^H(\bm{v})\bm{w}_k s_k + \sum_{j \neq k}^K \bm{h}_k^H(\bm{v})\bm{w}_j s_j + n_k
  \label{eq:signal_interference_decomposition}
\end{equation}
其中第一项为用户 $k$ 期望的有效信号分量，第二项为其他多用户间的共道多址干扰（Multi-User Interference, MUI），第三项为热噪声。

\subsection{瞬时信干噪比与和速率优化问题形式化}
根据式~\eqref{eq:signal_interference_decomposition}，用户 $k$ 处的瞬时接收信号与干扰加噪声比（SINR）可精确表征为：
\begin{equation}
  \gamma_k(\bm{W}, \bm{v}) = \frac{|\bm{h}_k^H(\bm{v})\bm{w}_k|^2}{\sum_{j \neq k}^K |\bm{h}_k^H(\bm{v})\bm{w}_j|^2 + \sigma_k^2}
  \label{eq:sinr_formulation}
\end{equation}
假设系统采用高斯信码，根据香农信息论，多用户下行信道的可达加权总和速率（Sum Rate, 单位：$\mathrm{bps/Hz}$）为：
\begin{equation}
  R_{\mathrm{sum}}(\bm{W}, \bm{v}) = \sum_{k=1}^K \log_2\left( 1 + \gamma_k(\bm{W}, \bm{v}) \right)
  \label{eq:sum_rate_def}
\end{equation}

本论文的核心目标是在满足基站发射功率限制与 RIS 各单元物理恒模相移约束下，联合优化基站有源预编码矩阵 $\bm{W}$ 与 RIS 无源反射相移向量 $\bm{v}$，以最大化系统下行总和速率。该主被动联合波束赋形问题形式化表述为：
\begin{subequations}
  \label{prob:joint_beamforming_p1}
  \begin{align}
    \max_{\bm{W}, \bm{v}} \quad & \sum_{k=1}^K \log_2\left( 1 + \frac{|\bm{h}_k^H(\bm{v})\bm{w}_k|^2}{\sum_{j \neq k}^K |\bm{h}_k^H(\bm{v})\bm{w}_j|^2 + \sigma_k^2} \right) \label{prob:p1_obj} \\
    \text{s.t.} \quad & \sum_{k=1}^K \|\bm{w}_k\|^2 \le P_{\max}, \label{prob:p1_power} \\
                      & |v_n| = 1, \quad \forall n = 1, 2, \dots, N \label{prob:p1_unimodular}
  \end{align}
\end{subequations}

优化问题~\eqref{prob:joint_beamforming_p1}呈现出极高的数学求解复杂性：一方面，目标函数~\eqref{prob:p1_obj}由多用户非凸对数分式求和构成，变量 $\bm{W}$ 与 $\bm{v}$ 在分子与分母中深度强耦合；另一方面，约束条件~\eqref{prob:p1_unimodular}施加了严格的单位圆单位模非凸等式约束。直接对其进行凸松弛往往会导致严重的维度灾难与次优极值陷阱。

% -*- coding: utf-8 -*-
\section{基于交替优化与分式规划的主被动波束协同设计}
\label{sec:beamforming_optimization}

针对多用户 MISO 和速率优化问题~\eqref{prob:joint_beamforming_p1}中发射预编码矩阵 $\bm{W}$ 与无源相移向量 $\bm{v}$ 的高度非凸非线性耦合，本节提出一套基于交替优化（Alternating Optimization, AO）与多比率分式规划（Fractional Programming, FP）的高效解耦协同优化算法。整体思路是：先用拉格朗日对偶变换把「对数」从目标函数中剥离，再用二次变换把「分式」化为二次型，从而使原本关于 $\bm{W}$ 与 $\bm{v}$ 均非凸的问题在两个子问题上分别获得凸结构或流形上的光滑结构。

\subsection{基于拉格朗日对偶与二次变换的多比率分式规划}
\label{sec:fp_transform}

由于目标函数中 SINR 处于对数函数内部且各项互为干扰，传统 Dinkelbach 标量分式规划无法直接处理多比率求和问题。为此，本文依次引入拉格朗日对偶变换与二次变换两级解耦。

\textbf{第一级：拉格朗日对偶变换。}为书写简洁，以下以自然对数表述目标函数；换算到以 2 为底仅相差常数因子 $1/\ln 2$，不影响最优解的位置。引入辅助变量 $\bm{\rho} = [\rho_1, \dots, \rho_K]^T \in \mathbb{R}_{+}^K$，利用恒等式
\begin{equation}
  \ln\left(1 + \gamma_k\right) = \max_{\rho_k > 0}\ \ln(1+\rho_k) - \rho_k + \frac{(1+\rho_k)\gamma_k}{1+\gamma_k},
  \label{eq:dual_transform}
\end{equation}
可将和速率中的每一项对数从 SINR 内部解耦。式~\eqref{eq:dual_transform}右端关于 $\rho_k$ 是严格凹的，令其导数为零即得驻点条件
\begin{equation}
  \frac{1}{1+\rho_k} - 1 + \frac{\gamma_k}{1+\gamma_k} = 0
  \quad\Longleftrightarrow\quad
  \rho_k^{\star} = \gamma_k(\bm{W}, \bm{v}),
  \label{eq:fp_update_rho}
\end{equation}
即辅助变量在最优处恰等于当前迭代点的 SINR。将 $\gamma_k$ 的定义~\eqref{eq:sinr_formulation}代入并整理，右端分式可写成
\begin{equation}
  \frac{(1+\rho_k)\gamma_k}{1+\gamma_k}
  = (1+\rho_k)\cdot\frac{|\bm{h}_k^H\bm{w}_k|^2}{\sum_{j=1}^{K}|\bm{h}_k^H\bm{w}_j|^2 + \sigma_k^2}
  = (1+\rho_k)\cdot\frac{|a_k|^2}{B_k + |a_k|^2},
  \label{eq:qt_transform}
\end{equation}
其中记 $a_k \triangleq \bm{h}_k^H(\bm{v})\bm{w}_k$ 为期望信号分量，$B_k \triangleq \sum_{j\neq k}|\bm{h}_k^H(\bm{v})\bm{w}_j|^2 + \sigma_k^2$ 为干扰加噪声功率。

\textbf{第二级：二次变换。}式~\eqref{eq:qt_transform}右端仍是分式。引入第二组辅助变量 $\mu_k$（$k = 1,\dots,K$），利用二次变换恒等式
\begin{equation}
  \frac{|a_k|^2}{B_k + |a_k|^2}
  = \max_{\mu_k \in \mathbb{C}}\ 2\mathrm{Re}\{\mu_k^* a_k\} - |\mu_k|^2\left(B_k + |a_k|^2\right),
  \label{eq:qt_identity}
\end{equation}
该式关于 $\mu_k$ 是严格凹的二次函数，最优解为
\begin{equation}
  \mu_k^{\star} = \frac{a_k}{B_k + |a_k|^2}
  = \frac{\bm{h}_k^H(\bm{v})\bm{w}_k}{\sum_{j=1}^{K}|\bm{h}_k^H(\bm{v})\bm{w}_j|^2 + \sigma_k^2}.
  \label{eq:fp_update_mu}
\end{equation}
将式~\eqref{eq:dual_transform}与式~\eqref{eq:qt_identity}逐项代入和速率，原问题~\eqref{prob:joint_beamforming_p1}等价重构为如下增广目标函数的极大化问题：
\begin{equation}
\begin{aligned}
  f_q(\bm{W}, \bm{v}, \bm{\rho}, \bm{\mu}) \triangleq \sum_{k=1}^K \Bigl[ &\ln(1 + \rho_k) - \rho_k \\
  &+ (1+\rho_k)\Bigl( 2\mathrm{Re}\{\mu_k^* \bm{h}_k^H(\bm{v})\bm{w}_k\} \\
  &\qquad - |\mu_k|^2 \Bigl( \sum_{j=1}^K |\bm{h}_k^H(\bm{v})\bm{w}_j|^2 + \sigma_k^2 \Bigr) \Bigr) \Bigr]
\end{aligned}
  \label{eq:fp_augmented_objective}
\end{equation}

\begin{theorem}[等价性]
\label{thm:fp_equivalence}
对任意可行 $(\bm{W}, \bm{v})$，有
\begin{equation}
  \max_{\bm{\rho},\bm{\mu}} f_q(\bm{W}, \bm{v}, \bm{\rho}, \bm{\mu}) = \sum_{k=1}^K \ln\left(1 + \gamma_k(\bm{W}, \bm{v})\right),
  \label{eq:fp_equivalence}
\end{equation}
且最优辅助变量由式~\eqref{eq:fp_update_rho}与式~\eqref{eq:fp_update_mu}唯一给出。因此交替优化 $(\bm{W},\bm{v})$ 与 $(\bm{\rho},\bm{\mu})$ 不改变原问题的最优值。
\end{theorem}

定理~\ref{thm:fp_equivalence}的意义在于：原问题中对数、分式、耦合三重困难被转移为「辅助变量闭式更新 + 两个结构化子问题」。辅助变量更新不涉及任何搜索，计算量为 $O(KM)$，可忽略；真正的计算负担集中在两个子问题上。

\subsection{子问题一：固定相移求解基站有源预编码}
\label{sec:sub_w}

在固定 $\bm{v}$、$\bm{\rho}$ 与 $\bm{\mu}$ 的条件下，目标函数式~\eqref{eq:fp_augmented_objective}关于基站有源预编码矩阵 $\bm{W}$ 退化为一个标准的凸二次规划（QP）问题：
\begin{subequations}
  \label{prob:sub1_w_opt}
  \begin{align}
    \min_{\bm{W}} \quad & \sum_{j=1}^K \bm{w}_j^H \left( \sum_{k=1}^K (1+\rho_k)|\mu_k|^2 \bm{h}_k(\bm{v})\bm{h}_k^H(\bm{v}) \right) \bm{w}_j - 2\sum_{k=1}^K (1+\rho_k)\mathrm{Re}\left\{ \mu_k^* \bm{h}_k^H(\bm{v})\bm{w}_k \right\} \label{prob:sub1_obj} \\
    \text{s.t.} \quad & \sum_{k=1}^K \|\bm{w}_k\|^2 \le P_{\max} \label{prob:sub1_power}
  \end{align}
\end{subequations}

利用拉格朗日乘子法，引入对偶乘子 $\lambda \ge 0$ 对应发射功率约束~\eqref{prob:sub1_power}，构造拉格朗日函数 $\mathcal{L}(\bm{W}, \lambda)$。对任意用户 $k$ 的波束赋形向量 $\bm{w}_k$ 求复梯度并令其为零：
\begin{equation}
  \nabla_{\bm{w}_k^*} \mathcal{L} = \left( \sum_{j=1}^K (1+\rho_j)|\mu_j|^2 \bm{h}_j(\bm{v})\bm{h}_j^H(\bm{v}) + \lambda \bm{I}_M \right) \bm{w}_k - (1+\rho_k)\mu_k \bm{h}_k(\bm{v}) = \bm{0}
  \label{eq:w_k_derivative}
\end{equation}
由此可得最优发射预编码向量 $\bm{w}_k^*$ 的显式解析闭式解：
\begin{equation}
  \bm{w}_k^*(\lambda) = \left( \bm{A}(\bm{v}, \bm{\mu}) + \lambda \bm{I}_M \right)^{-1} (1+\rho_k)\mu_k \bm{h}_k(\bm{v})
  \label{eq:optimal_wk_closed_form}
\end{equation}
其中 $\bm{A}(\bm{v}, \bm{\mu}) \triangleq \sum_{j=1}^K (1+\rho_j)|\mu_j|^2 \bm{h}_j(\bm{v})\bm{h}_j^H(\bm{v}) \in \mathbb{C}^{M \times M}$ 为半正定等效协方差矩阵。对偶乘子 $\lambda \ge 0$ 可通过一维二分搜索法（Bisection Method）在区间 $[0, \lambda_{\max}]$ 内精确确定，满足互补松弛条件：
\begin{equation}
  \lambda \left( \sum_{k=1}^K \|\bm{w}_k^*(\lambda)\|^2 - P_{\max} \right) = 0
  \label{eq:slackness_condition}
\end{equation}
由于 $\sum_k\|\bm{w}_k^*(\lambda)\|^2$ 关于 $\lambda$ 单调递减，二分搜索的收敛性有保证。该子问题因此具有「闭式解 + 一维搜索」的结构，单次求解复杂度仅为 $O(M^3 + KM^2 + M^2\log(1/\epsilon))$。

需要指出，子问题~\eqref{prob:sub1_w_opt}的凸性依赖于辅助变量 $(\bm{\rho},\bm{\mu})$ 被固定。若直接从原始和速率出发优化 $\bm{W}$，问题在 $\bm{W}$ 上并不凸；正是分式规划把非凸性「搬进」了辅助变量，才使 $\bm{W}$ 子问题成为标准 QP。这是本文采用 FP 而非直接梯度法的根本原因。

\subsection{子问题二：固定预编码求解 RIS 无源相移向量}
\label{sec:sub_v}

在给定基站预编码矩阵 $\bm{W}$ 及辅助变量 $\bm{\rho}, \bm{\mu}$ 的条件下，将复合信道 $\bm{h}_k^H(\bm{v}) = \bm{h}_{d,k}^H + \bm{v}^H \bm{H}_{\mathrm{casc},k}$ 代入目标函数式~\eqref{eq:fp_augmented_objective}中展开。注意到 $\bm{h}_k^H(\bm{v})$ 关于 $\bm{v}$ 是仿射的，因此 $|\bm{h}_k^H\bm{w}_j|^2$ 是 $\bm{v}$ 的二次型，$\mathrm{Re}\{\mu_k^*\bm{h}_k^H\bm{w}_k\}$ 是 $\bm{v}$ 的线性项。

经代数重组，关于相移向量 $\bm{v} \in \mathbb{C}^{N \times 1}$ 的优化问题可等价整理为标准的复二次型极小化问题：
\begin{subequations}
  \label{prob:sub2_v_opt}
  \begin{align}
    \min_{\bm{v}} \quad & \bm{v}^H \bm{\Psi} \bm{v} - 2\mathrm{Re}\{\bm{v}^H \bm{q}\} \label{prob:sub2_obj} \\
    \text{s.t.} \quad & |v_n| = 1, \quad \forall n = 1, 2, \dots, N \label{prob:sub2_unimodular}
  \end{align}
\end{subequations}
其中矩阵 $\bm{\Psi} \in \mathbb{C}^{N \times N}$ 与向量 $\bm{q} \in \mathbb{C}^{N \times 1}$ 分别定义为：
\begin{align}
  \bm{\Psi} &\triangleq \sum_{k=1}^K (1+\rho_k)|\mu_k|^2 \sum_{j=1}^K \bm{H}_{\mathrm{casc},k} \bm{w}_j \bm{w}_j^H \bm{H}_{\mathrm{casc},k}^H \label{eq:psi_matrix_def} \\
  \bm{q} &\triangleq \sum_{k=1}^K (1+\rho_k)\left[ \mu_k \bm{H}_{\mathrm{casc},k}\bm{w}_k - |\mu_k|^2 \sum_{j=1}^K \bm{H}_{\mathrm{casc},k} \bm{w}_j \bm{w}_j^H \bm{h}_{d,k} \right] \label{eq:q_vector_def}
\end{align}
根据定义，矩阵 $\bm{\Psi}$ 为半正定厄米特矩阵（$\bm{\Psi} \succeq \bm{0}$），因此式~\eqref{prob:sub2_obj}在无约束时是凸的；非凸性完全来自单位模约束~\eqref{prob:sub2_unimodular}。

由于约束~\eqref{prob:sub2_unimodular}构成了由 $N$ 个复平面单位圆构成的黎曼流形 $\mathcal{M} = \{\bm{v} \in \mathbb{C}^N : |v_n|=1, \forall n\}$，本节采用复流形共轭梯度法（Riemannian Conjugate Gradient, RCG）进行直接求解。算法~\ref{alg:rcg}给出完整流程。

\begin{algorithm}[t]
\caption{单位模约束下的黎曼共轭梯度法}
\label{alg:rcg}
\begin{algorithmic}[1]
\REQUIRE $\bm{\Psi} \succeq \bm{0}$, $\bm{q}$, 初始点 $\bm{v}_0 \in \mathcal{M}$
\ENSURE 局部最优相移向量 $\bm{v}$
\STATE 初始化 $\bm{d}_0 = -\mathrm{grad}\,g(\bm{v}_0)$，$t \leftarrow 0$
\REPEAT
  \STATE 计算欧氏梯度 $\nabla g(\bm{v}_t) = 2\bm{\Psi}\bm{v}_t - 2\bm{q}$
  \STATE 投影至切空间：$\mathrm{grad}\,g(\bm{v}_t) = \nabla g(\bm{v}_t) - \mathrm{Re}\{\nabla g(\bm{v}_t)\circ\bm{v}_t^*\}\circ\bm{v}_t$
  \STATE Armijo 线搜索确定步长 $\alpha_t$，使目标函数充分下降
  \STATE 收缩回流形：$\bm{v}_{t+1} = \mathcal{R}_{\bm{v}_t}(\alpha_t \bm{d}_t)$
  \STATE 计算 Polak-Ribière 系数 $\beta_t = \frac{\langle \mathrm{grad}\,g(\bm{v}_{t+1}),\ \mathrm{grad}\,g(\bm{v}_{t+1}) - \mathrm{grad}\,g(\bm{v}_t)\rangle}{\langle \mathrm{grad}\,g(\bm{v}_t),\ \mathrm{grad}\,g(\bm{v}_t)\rangle}$
  \STATE 更新搜索方向 $\bm{d}_{t+1} = -\mathrm{grad}\,g(\bm{v}_{t+1}) + \beta_t \mathcal{T}(\bm{d}_t)$
  \STATE $t \leftarrow t+1$
\UNTIL{$\|\mathrm{grad}\,g(\bm{v}_t)\| \le \epsilon$ 或达到最大迭代次数}
\RETURN $\bm{v}_t$
\end{algorithmic}
\end{algorithm}

其中两个关键算子的具体形式为：
\begin{enumerate}
  \item \textbf{黎曼梯度投影}：计算欧几里得空间梯度 $\nabla_{\bm{v}} g(\bm{v}) = 2\bm{\Psi}\bm{v} - 2\bm{q}$，并将其切空间投影至复流形 $\mathcal{M}$ 的切空间 $T_{\bm{v}}\mathcal{M}$：
    \begin{equation}
      \mathrm{grad} \, g(\bm{v}) = \nabla_{\bm{v}} g(\bm{v}) - \mathrm{Re}\left\{ \nabla_{\bm{v}} g(\bm{v}) \circ \bm{v}^* \right\} \circ \bm{v}
      \label{eq:riemannian_grad}
    \end{equation}
    其中 $\circ$ 表示 Hadamard 逐元素积。该投影的几何含义是剔除梯度中沿半径方向（即破坏单位模）的分量。
  \item \textbf{流形回退收缩 (Retraction)}：沿着共轭搜索方向 $\bm{d}$ 进行 Armijo 线搜索步长更新后，通过如下收缩算子将迭代点投影回单位模复圆流形上：
    \begin{equation}
      \mathcal{R}_{\bm{v}}(\alpha \bm{d}) = \left[ \frac{(\bm{v} + \alpha \bm{d})_1}{|(\bm{v} + \alpha \bm{d})_1|}, \dots, \frac{(\bm{v} + \alpha \bm{d})_N}{|(\bm{v} + \alpha \bm{d})_N|} \right]^T
      \label{eq:manifold_retraction}
    \end{equation}
\end{enumerate}

流形共轭梯度法避开了 SDR 升维将变量转化为 $N \times N$ 矩阵的巨大计算开销，直接在 $N$ 维流形上实现超线性收敛。相较于仅沿负梯度方向的黎曼梯度下降，共轭方向利用了历史梯度信息，在病态二次型（$\bm{\Psi}$ 条件数较大）上能显著减少迭代次数；这一点在第四章的收敛性实验中表现为：共轭梯度在十余次迭代内即可把和速率提升到稳定值的 99\% 以上。

\subsection{交替优化框架与复杂度汇总}
\label{sec:ao_summary}

将上述两块拼合，得到完整的交替优化流程：固定 $(\bm{\rho},\bm{\mu})$ 更新 $\bm{W}$（闭式 + 二分），固定 $\bm{W}$ 更新 $\bm{v}$（RCG），随后用式~\eqref{eq:fp_update_rho}与式~\eqref{eq:fp_update_mu}刷新辅助变量。由于每一步都在增广目标函数 $f_q$ 上做精确或单调提升的极大化，而 $f_q$ 有上界，整个交替过程单调不减且必然收敛到 $f_q$ 的驻点。需要强调的是，由于原问题非凸，所得驻点是局部最优而非全局最优，其质量依赖初始点；第四章将给出多初值重启策略以缓解这一问题。

单次交替迭代的复杂度为 $O(M^3 + KM^2 + N^2 M + N^3)$，其中 $M^3$ 来自 $\bm{W}$ 子问题的矩阵求逆（可用 Sherman-Morrison 加速至 $O(KM^2)$），$N^3$ 来自 $\bm{v}$ 子问题中 $\bm{\Psi}$ 的构造与矩阵—向量乘。在 $M=8$、$N=128$、$K=4$ 的典型配置下，单次迭代的浮点运算量约为 $4\times 10^{6}$ 量级，可在毫秒级完成，满足实时重构需求。

% -*- coding: utf-8 -*-
\section{基于正交 DFT 码本的低开销级联信道估计协议}
\label{sec:channel_estimation}

前述主被动协同波束赋形算法的最优性严格依赖于基站端获取精确的信道状态信息（CSI）。然而，由于无源 RIS 不具备基带解码与射频处理单元，基站与用户之间无法分别直接测量独立的“基站-表面”信道 $\bm{G}$ 与“表面-用户”信道 $\bm{h}_{r,k}^H$。收发机仅能观测到两者的级联响应矩阵 $\bm{H}_{\mathrm{casc},k} = \mathrm{diag}(\bm{h}_{r,k}^H)\bm{G} \in \mathbb{C}^{N \times M}$ 以及直连信道 $\bm{h}_{d,k}^H \in \mathbb{C}^{1 \times M}$。本节设计一种基于时域正交导频与离散傅里叶变换（DFT）反射码本的低开销信道估计协议。

\subsection{上行导频传输与级联信道矩阵向量化}
基于时分双工（TDD）系统的信道互易性定理，下行级联 CSI 可通过上行导频训练进行高效获取。假定在相干时隙的训练阶段，各用户并发发送相互正交的导频序列 $\bm{S}_p \in \mathbb{C}^{K \times \tau_p}$，满足：
\begin{equation}
  \bm{S}_p \bm{S}_p^H = \tau_p \bm{I}_K
  \label{eq:pilot_orthogonality}
\end{equation}
其中 $\tau_p \ge K$ 为单次导频块符号长度。基站通过乘以各用户的匹配滤波矩阵 $\bm{S}_p^H$，可完全无失真解耦不同用户间的导频响应。对于用户 $k$，在第 $p$ 个导频子时隙（对应 RIS 反射相移配置为 $\bm{v}_p \in \mathbb{C}^{N \times 1}$）下，基站接收的解耦后上行信号矩阵 $\bm{Y}_{k,p} \in \mathbb{C}^{M \times 1}$ 满足：
\begin{equation}
  \bm{y}_{k,p} = \bm{h}_{d,k} + \bm{G}^H \mathrm{diag}(\bm{v}_p) \bm{h}_{r,k} + \bar{\bm{n}}_{k,p}
  \label{eq:uplink_pilot_received}
\end{equation}
其中 $\bar{\bm{n}}_{k,p} \sim \mathcal{CN}(\bm{0}, \frac{\sigma^2}{\tau_p}\bm{I}_M)$ 为经过匹配滤波后的等效高斯白噪声。

利用代数恒等式 $\mathrm{diag}(\bm{v}_p)\bm{h}_{r,k} = \mathrm{diag}(\bm{h}_{r,k})\bm{v}_p$，式~\eqref{eq:uplink_pilot_received}中级联链路项可重写为：
\begin{equation}
  \bm{G}^H \mathrm{diag}(\bm{v}_p) \bm{h}_{r,k} = \bm{H}_{\mathrm{casc},k}^H \bm{v}_p = \left( \bm{v}_p^T \otimes \bm{I}_M \right) \mathrm{vec}(\bm{H}_{\mathrm{casc},k}^H)
  \label{eq:kronecker_vectorization}
\end{equation}
其中 $\otimes$ 表示 Kronecker 积，$\mathrm{vec}(\cdot)$ 表示矩阵按列向量化算子。

定义全级联信道聚合向量为：
\begin{equation}
  \bar{\bm{h}}_k \triangleq \begin{bmatrix} \bm{h}_{d,k} \\ \mathrm{vec}(\bm{H}_{\mathrm{casc},k}^H) \end{bmatrix} \in \mathbb{C}^{(N+1)M \times 1}
  \label{eq:aggregated_channel_vector}
\end{equation}
在 $P$ 个连续训练子时隙内，RIS 依次变换 $P$ 组不同的反射相移图案 $\bm{v}_1, \bm{v}_2, \dots, \bm{v}_P$。构建 RIS 导频反射设计矩阵：
\begin{equation}
  \bm{\Phi} \triangleq \begin{bmatrix}
    1 & \bm{v}_1^T \\
    1 & \bm{v}_2^T \\
    \vdots & \vdots \\
    1 & \bm{v}_P^T
  \end{bmatrix} \in \mathbb{C}^{P \times (N+1)}
  \label{eq:ris_pilot_matrix_phi}
\end{equation}
将所有 $P$ 个时隙的基站接收信号级联堆叠，可得紧凑的观测线性模型：
\begin{equation}
  \bar{\bm{y}}_k \triangleq \begin{bmatrix} \bm{y}_{k,1} \\ \bm{y}_{k,2} \\ \vdots \\ \bm{y}_{k,P} \end{bmatrix} = \left( \bm{\Phi} \otimes \bm{I}_M \right) \bar{\bm{h}}_k + \bar{\bm{n}}_k
  \label{eq:compact_observation_model}
\end{equation}
其中 $\bar{\bm{n}}_k \in \mathbb{C}^{PM \times 1}$ 为综合噪声向量，协方差矩阵为 $\bm{R}_n = \frac{\sigma^2}{\tau_p}\bm{I}_{PM}$。

\subsection{基于正交 DFT 反射码本的最小二乘解耦估计}
由观测方程~\eqref{eq:compact_observation_model}可知，要保证未知信道向量 $\bar{\bm{h}}_k$ 具有唯一确定解，观测矩阵必须满足列满秩条件，即必需且仅需训练子时隙数 $P \ge N+1$。

为了最小化噪声放大效应并达到最小估计均方误差（MSE），RIS 导频反射矩阵 $\bm{\Phi}$ 的列向量之间必须保持严格的正交性。为此，本节构造基于截断离散傅里叶变换（DFT）的正交码本：
令 $\bm{F}_{P} \in \mathbb{C}^{P \times P}$ 表示标准 $P$ 阶归一化 DFT 矩阵，其第 $(p, q)$ 项元素定义为：
\begin{equation}
  [\bm{F}_P]_{p, q} = \frac{1}{\sqrt{P}} \exp\left( -j \frac{2\pi (p-1)(q-1)}{P} \right), \quad 1 \le p, q \le P
  \label{eq:dft_matrix_element}
\end{equation}
选取 $P = N+1$（或当 $N+1$ 非 2 的幂次时选取 $P = 2^{\lceil \log_2(N+1) \rceil}$），直接令导频设计矩阵 $\bm{\Phi}$ 由 $\bm{F}_P$ 的前 $N+1$ 列构成：
\begin{equation}
  \bm{\Phi} = \sqrt{P} [\bm{F}_P]_{:, 1:(N+1)}
  \label{eq:optimal_dft_codebook}
\end{equation}
显然，该设计天然满足行恒模物理条件 $|[\bm{\Phi}]_{p, n+1}| = 1$，且满足完美列正交性：
\begin{equation}
  \bm{\Phi}^H \bm{\Phi} = P \bm{I}_{N+1}
  \label{eq:phi_column_orthogonality}
\end{equation}

在式~\eqref{eq:phi_column_orthogonality}的列正交性保障下，基站端可采用计算复杂度极低的最小二乘（LS）解析估计器直接闭式恢复全信道：
\begin{equation}
  \hat{\bar{\bm{h}}}_k^{\mathrm{LS}} = \left( (\bm{\Phi}^H \bm{\Phi})^{-1}\bm{\Phi}^H \otimes \bm{I}_M \right) \bar{\bm{y}}_k = \frac{1}{P} \left( \bm{\Phi}^H \otimes \bm{I}_M \right) \bar{\bm{y}}_k
  \label{eq:ls_closed_form_estimator}
\end{equation}
计算过程完全避免了高维矩阵求逆运算，单用户信道重构复杂度仅为 $\mathcal{O}(PM)$。
从重构向量 $\hat{\bar{\bm{h}}}_k^{\mathrm{LS}}$ 中，前 $M$ 个分量直接提取出基站至用户 $k$ 的直连信道估计 $\hat{\bm{h}}_{d,k}$；后 $NM$ 个分量经逆向量化反折叠，即刻精确恢复出级联信道矩阵 $\hat{\bm{H}}_{\mathrm{casc},k}$。
整个估计协议仅需 $P = N+1$ 个时隙即可同时完成直连链路与级联反射链路的完全正交解耦估计，导频开销达到理论下界。

% -*- coding: utf-8 -*-
\section{算法收敛性分析、CRLB推导与计算复杂度}
\label{sec:convergence_and_complexity}

为了严谨确立本章提出的主被动协同波束赋形与信道估计算法在数学与信息论层面的可靠性，本节给出交替优化算法的单调收敛性证明、信道估计的克拉美罗下界（CRLB）推导，以及计算复杂度的渐近阶次分析。

\subsection{交替优化波束赋形算法的单调收敛性证明}

\begin{theorem}[交替优化算法单调收敛性]
  \label{thm:ao_convergence}
  由式~\eqref{eq:fp_update_rho}、\eqref{eq:fp_update_mu}、\eqref{eq:optimal_wk_closed_form}与流形更新式~\eqref{eq:manifold_retraction}构成的迭代序列 $\{ \bm{W}^{(t)}, \bm{v}^{(t)} \}$，其对应的原始系统和速率目标函数值序列 $\{ R_{\mathrm{sum}}(\bm{W}^{(t)}, \bm{v}^{(t)}) \}$ 随外层迭代次数 $t$ 单调非递减，且序列严格收敛至非凸和速率优化问题~\eqref{prob:joint_beamforming_p1}的卡罗什-库恩-塔克（KKT）一阶驻点。
\end{theorem}

\begin{proof}
  令 $f_q(\bm{W}, \bm{v}, \bm{\rho}, \bm{\mu})$ 表示由式~\eqref{eq:fp_augmented_objective}定义的分式规划等价增广目标函数。根据二次变换等价性定理，对于任意给定的 $(\bm{W}^{(t)}, \bm{v}^{(t)})$，当辅助变量更新为最优值 $\bm{\rho}^{(t+1)}$ 与 $\bm{\mu}^{(t+1)}$ 时，增广函数与原目标函数严格相等：
  \begin{equation}
    R_{\mathrm{sum}}(\bm{W}^{(t)}, \bm{v}^{(t)}) = f_q\left( \bm{W}^{(t)}, \bm{v}^{(t)}, \bm{\rho}^{(t+1)}, \bm{\mu}^{(t+1)} \right)
    \label{eq:proof_step1}
  \end{equation}
  在第一步优化中，固定 $(\bm{v}^{(t)}, \bm{\rho}^{(t+1)}, \bm{\mu}^{(t+1)})$，由于解 $\bm{W}^{(t+1)}$ 为凸二次规划子问题~\eqref{prob:sub1_w_opt}的全局最优闭式解，因此有：
  \begin{equation}
    f_q\left( \bm{W}^{(t)}, \bm{v}^{(t)}, \bm{\rho}^{(t+1)}, \bm{\mu}^{(t+1)} \right) \le f_q\left( \bm{W}^{(t+1)}, \bm{v}^{(t)}, \bm{\rho}^{(t+1)}, \bm{\mu}^{(t+1)} \right)
    \label{eq:proof_step2}
  \end{equation}
  在第二步优化中，固定 $(\bm{W}^{(t+1)}, \bm{\rho}^{(t+1)}, \bm{\mu}^{(t+1)})$，黎曼流形共轭梯度法（RCG）沿切空间黎曼负梯度方向并配合 Armijo 步长准则进行流形回退收缩更新，保证子问题~\eqref{prob:sub2_v_opt}的目标函数单调递减（即对应原极大化目标函数非递减）：
  \begin{equation}
    f_q\left( \bm{W}^{(t+1)}, \bm{v}^{(t)}, \bm{\rho}^{(t+1)}, \bm{\mu}^{(t+1)} \right) \le f_q\left( \bm{W}^{(t+1)}, \bm{v}^{(t+1)}, \bm{\rho}^{(t+1)}, \bm{\mu}^{(t+1)} \right)
    \label{eq:proof_step3}
  \end{equation}
  根据式~\eqref{eq:fp_augmented_objective}的数学结构，对于任意给定的可行解 $(\bm{W}^{(t+1)}, \bm{v}^{(t+1)})$，原始目标函数值恒为增广目标函数关于辅助变量 $(\bm{\rho}, \bm{\mu})$ 的上确界：
  \begin{equation}
    f_q\left( \bm{W}^{(t+1)}, \bm{v}^{(t+1)}, \bm{\rho}^{(t+1)}, \bm{\mu}^{(t+1)} \right) \le \max_{\bm{\rho}, \bm{\mu}} f_q\left( \bm{W}^{(t+1)}, \bm{v}^{(t+1)}, \bm{\rho}, \bm{\mu} \right) = R_{\mathrm{sum}}(\bm{W}^{(t+1)}, \bm{v}^{(t+1)})
    \label{eq:proof_step4}
  \end{equation}
  综合式~\eqref{eq:proof_step1}至式~\eqref{eq:proof_step4}，可得严格的不等式链：
  \begin{equation}
    R_{\mathrm{sum}}(\bm{W}^{(t)}, \bm{v}^{(t)}) \le R_{\mathrm{sum}}(\bm{W}^{(t+1)}, \bm{v}^{(t+1)})
    \label{eq:proof_chain}
  \end{equation}
  此外，在满足最大总发射功率约束 $\sum_{k=1}^K \|\bm{w}_k\|^2 \le P_{\max}$ 与有限信道增益条件下，多用户总和速率 $R_{\mathrm{sum}}(\bm{W}, \bm{v})$ 必然存在有限的数学物理上界：
  \begin{equation}
    R_{\mathrm{sum}}(\bm{W}, \bm{v}) \le K \log_2\left( 1 + \frac{P_{\max} \max_k \|\bm{h}_k^H(\bm{v})\|^2}{\sigma^2} \right) < \infty
    \label{eq:sum_rate_upper_bound}
  \end{equation}
  由单调有界收敛定理（Monotone Convergence Theorem），目标函数值序列必然单调收敛至有限极限值，定理证毕。
\end{proof}

\subsection{DFT 导频信道估计误差的克拉美罗下界 (CRLB) 推导}

为了评估基于正交 DFT 码本的最小二乘（LS）信道估计器的渐近有效性，本小节推导未知信道向量 $\bar{\bm{h}}_k$ 的无偏估计克拉美罗下界（CRLB）。

\begin{theorem}[级联信道估计 CRLB]
  \label{thm:crlb_derivation}
  在复加性高斯白噪声信道环境下，用户 $k$ 的全级联信道向量 $\bar{\bm{h}}_k \in \mathbb{C}^{(N+1)M \times 1}$ 的任意无偏估计量 $\hat{\bar{\bm{h}}}_k$ 的协方差矩阵满足下界：
  \begin{equation}
    \mathrm{Cov}\left( \hat{\bar{\bm{h}}}_k \right) \succeq \bm{J}^{-1} = \frac{\sigma^2}{P \tau_p} \bm{I}_{(N+1)M}
    \label{eq:crlb_matrix_bound}
  \end{equation}
  且由式~\eqref{eq:ls_closed_form_estimator}给出的正交 DFT 最小二乘估计器达到该理论下界，其均方误差为：
  \begin{equation}
    \mathrm{MSE}(\hat{\bar{\bm{h}}}_k^{\mathrm{LS}}) = \mathrm{Tr}(\bm{J}^{-1}) = \frac{(N+1)M \sigma^2}{P \tau_p}
    \label{eq:crlb_trace_mse}
  \end{equation}
\end{theorem}

\begin{proof}
  根据观测模型式~\eqref{eq:compact_observation_model}，接收信号向量 $\bar{\bm{y}}_k \sim \mathcal{CN}\left( (\bm{\Phi} \otimes \bm{I}_M)\bar{\bm{h}}_k, \frac{\sigma^2}{\tau_p}\bm{I}_{PM} \right)$。对数似然函数关于未知复信道向量的 Fisher 信息矩阵（FIM）$\bm{J}$ 可严格计算为：
  \begin{equation}
    \bm{J} = \left( \bm{\Phi} \otimes \bm{I}_M \right)^H \left( \frac{\sigma^2}{\tau_p}\bm{I}_{PM} \right)^{-1} \left( \bm{\Phi} \otimes \bm{I}_M \right) = \frac{\tau_p}{\sigma^2} \left( \bm{\Phi}^H \bm{\Phi} \otimes \bm{I}_M \right)
    \label{eq:fim_derivation}
  \end{equation}
  将最优 DFT 码本的正交性质式~\eqref{eq:phi_column_orthogonality}（即 $\bm{\Phi}^H \bm{\Phi} = P \bm{I}_{N+1}$）代入式~\eqref{eq:fim_derivation}：
  \begin{equation}
    \bm{J} = \frac{\tau_p}{\sigma^2} \left( P \bm{I}_{N+1} \otimes \bm{I}_M \right) = \frac{P \tau_p}{\sigma^2} \bm{I}_{(N+1)M}
    \label{eq:fim_orthogonal_result}
  \end{equation}
  对对角矩阵 $\bm{J}$ 直接求逆，即得 CRLB 协方差下界~\eqref{eq:crlb_matrix_bound}。对于线性高斯模型中的 LS 估计器，其误差协方差矩阵为：
  \begin{equation}
    \mathrm{Cov}\left( \hat{\bar{\bm{h}}}_k^{\mathrm{LS}} \right) = \left( (\bm{\Phi} \otimes \bm{I}_M)^H \bm{R}_n^{-1} (\bm{\Phi} \otimes \bm{I}_M) \right)^{-1} = \bm{J}^{-1}
    \label{eq:ls_attains_crlb}
  \end{equation}
  即 LS 估计器的协方差矩阵恒等达到 CRLB 下界，证明了正交 DFT 码本设计的最小方差无偏性（MVUE），定理证毕。
\end{proof}

值得指出的是，在最小导频开销配置 $P = N+1$ 条件下，代入式~\eqref{eq:crlb_trace_mse}可得单个信道系数的平均估计误差为 $\frac{\sigma^2}{(N+1)\tau_p}$，完全消除了由于反射微元增多带来的误差扩散。

\subsection{算法计算复杂度对比分析}
本节算法的整体计算复杂度由基站波束子问题求解、RIS 相移流形优化以及信道估计三部分组成：
\begin{enumerate}
  \item \textbf{基站有源预编码求解}：构建矩阵 $\bm{A}(\bm{v}, \bm{\mu})$ 需要 $\mathcal{O}(KM^2)$ 次浮点运算；进行矩阵求逆需要 $\mathcal{O}(M^3)$；二分搜索迭代通常在 $I_{\mathrm{bis}} \le 15$ 次内收敛。子问题一单次更新复杂度为 $\mathcal{O}(KM^2 + M^3)$。
  \item \textbf{RIS 无源相移求解}：构造矩阵 $\bm{\Psi}$ 与向量 $\bm{q}$ 的计算量为 $\mathcal{O}(K(N^2M + NM^2))$；复流形共轭梯度法（RCG）单次内层迭代仅涉及矩阵-向量乘法，复杂度为 $\mathcal{O}(N^2)$；若平均经过 $I_{\mathrm{RCG}} \approx 20\sim 30$ 次内层迭代收敛，子问题二总复杂度为 $\mathcal{O}(K(N^2M + NM^2) + I_{\mathrm{RCG}} N^2)$。
  \item \textbf{DFT 级联信道估计}：得益于 $\bm{\Phi}$ 的正交性，LS 估计器仅需执行矩阵相乘式~\eqref{eq:ls_closed_form_estimator}，无需任何矩阵求逆，单用户估计复杂度仅为 $\mathcal{O}(PM) = \mathcal{O}((N+1)M)$。
\end{enumerate}

与现有文献中基于半正定松弛（SDR）的算法进行复杂度横向对比：SDR 将相移向量松弛为 $N \times N$ 的半正定矩阵，采用内点法求解的计算复杂度高达 $\mathcal{O}(N^{4.5})$；当 RIS 单元数 $N=128$ 时，SDR 单次迭代浮点运算次数超过 $10^9$ 量级，而本论文所提出的 RCG 算法单次迭代计算量仅为 $\mathcal{O}(N^2) \approx 1.6 \times 10^4$，计算效率提升了 5 个数量级，完全满足 6G 毫秒级相干块内的实时调度需求。


\section{本章小结}
本章完成了面向 6G 的 RIS 辅助多用户通信系统的完整理论体系推导与算法设计：
\begin{enumerate}
  \item 形式化建立了包含基站直连链路与无源级联反射链路的复合信道传输模型，推导了严格的多用户 SINR 与和速率闭式表达式，并揭示了级联矩阵「两段信道逐元素相乘」的秩一哈达玛调制结构；
  \item 提出了基于二次变换分式规划与黎曼流形共轭梯度的交替优化（AO）算法，将复杂的双变量非凸约束问题解耦为两个高效闭式/拟闭式可解凸子问题，并说明了为何放弃升维松弛（SDR）路线——在 $N = 256$ 时其协方差变量规模达百万量级；
  \item 设计了基于正交 DFT 反射矩阵的低开销级联信道估计协议，将导频开销压至 $P = N+1$ 的理论下界，推导并给出了其估计误差的克拉美罗下界（CRLB），同时论证了在 $T < N+1$ 时压缩感知先验并非天然适用；
  \item 证明了交替优化算法的目标函数单调非递减性与收敛性，给出了单次迭代与整体收敛的计算复杂度界限，并通过与 SDR、SCA、ADMM 的复杂度横向对比说明了本方法的可实现性。
\end{enumerate}

本章的结论也存在明确的适用边界，需在后续章节中如实检验。其一，收敛性定理保证的是「收敛到驻点」而非全局最优，最终解依赖初值，因此第四章的收敛性实验只报告单调性与收敛速度，不宣称最优性。其二，CRLB 的成立以「训练期内信道静止、噪声独立同分布复高斯」为前提，实际系统的相位噪声与时变性会使实测误差高于下界，第四章的蒙特卡洛比值因此允许小幅偏离 1。其三，本章未实现 SDR 数值基准，复杂度对比仅停留在渐近阶次与单次迭代量级层面，这一取舍已在 3.4.3 节明确说明。

本章所推导的算法机制与理论界限，为第四章系统级仿真与数值实验提供了坚实完备的方法论依据。
